\section{Background: Information Extraction in the Legal Domain}


\todo{from claw}
\todo{add also from wise for investigations}

\todo{following is from aixia}
\subsection{NEEL for Legal Documents}

Most of the previous work on NEEL for legal documents has focused on the NER task only. The first NER approaches are based on handcrafted rules and statistical models~\cite{ccetindaug2023name_legal_turkish}, such as conditional random fields~(CRFs)~\cite{Lafferty2001}. More recent approaches to NER started using BiLSTM-based models combined with CRFs for Brazilian and German legal texts~\cite{ccetindaug2023name_legal_turkish}.
After the advent of transformers~\cite{vaswani_transformers_2017} and BERT~\cite{devlin-etal-2019-bert}, which obtained impressive performance in multiple NLP tasks, LEGAL-BERT, specialized in the legal domain, has been released~\cite{chalkidis-etal-2020-legal}. Later, some work compared LEGAL-BERT with previous approaches~\cite{Keshavarz2022} for NER finding that LEGAL-BERT performance is comparable to simple models (LSTMs, CNNs). 

A few approaches have studied NEL in legal texts. One of the first approaches applied NER and NEL on a corpus of judgments of the European Court of Human Rights~\cite{cardellino_etal_2017_lowcost}. As the background KB, they use a legal-specific ontology enriched with YAGO~\cite{mahdisoltani2013yago3,yago2_2013,yago2s_demo_2013,yago_core_2007} after an alignment procedure. Another approach targets the NEL task only on the EUR-Lex law article dataset~\cite{Elnaggar2018}. Their NEL system is trained using transfer learning. We study the end-to-end combination of NER, NEL, and NIL prediction, and we use a more recent NEL approach~\cite{blink} trained on a large Italian Wikipedia corpus without fine-tuning on court judgment data. Another study combined BERT~\cite{devlin-etal-2019-bert} with rule-based techniques for NER and coupled it with an off-the-shelves NEL service to extract entities from court decisions in the Finnish language~\cite{Tamper2020}; NEL is performed with a popularity-based approach. This study is the most similar to ours; however, we focus on the NIL prediction problem, and we use a BERT-based NEL approach; also, in this paper, we discuss only the performance of neural algorithms. 

Some work that studies NEL with NIL prediction is also evaluated on documents that are related to the legal domain~\cite{klie-etal-2020-zero} (the depositions of the 1641 Irish rebellion.%~\cite{tcd1641depositions}\todo{magari rimuovi citazione}.
However, to the best of our knowledge, no prior work has investigated end-to-end NEEL considering the NIL prediction problem on recent legal data. 

\subsection{NEEL in the Italian Language}
% \subsubsection{Datasets}
Italian datasets for NER include multilingual resources~\cite{tedeschi-navigli-2022-multinerd,NOTHMAN2013151}, and domain-specific datasets, such as \cite{CATELLI2020106779} in the medical domain.
Similarly, Italian NEL datasets comprise multilingual ones, i.e. VoxEL~\cite{rosales2018voxel} and resources based on micro-posts~\cite{basile2016neelit}.

% \subsubsection{Enity Extraction in Italian}
Among the ready-to-use NEEL libraries for the Italian language, notable options are SpaCy~\cite{honnibal2020spacy} and Tint~\cite{2018tint2} for NER and DBpedia Spotlight~\cite{mendes_2011_spotlight,Daiber2013spotlight} for both NER and NEL.\todo{add dbpedia spotlight as library for nel in background}
SpaCy provides pre-trained NER models of different sizes (small, medium and large), but currently does not provide any pre-trained transformer-based model for Italian. Tint performs NER with a combination of CRFs taggers and rule-based systems for dates and money.
DBpedia Spotlight is a ready-to-use tool that recognizes and links entity mentions to DBpedia~\cite{dbpedia}.
